\chapter{Σύστημα, Δεδομένα και Μεθοδολογία}

\section{IEEE 39 και Ambient Scenario}

Για την παραγωγή των \textit{ambient} δεδομένων εκτελέστηκε δυναμική προσομοίωση του \textit{IEEE 39-Bus System}~\cite{digsilent39} στο περιβάλλον \textit{PowerFactory}. Συγκεκριμένα, εφαρμόστηκε στοχαστική διέγερση σε όλα τα φορτία του συστήματος και για κάθε φορτίο παράχθηκαν δύο ανεξάρτητες χρονοσειρές θορύβου, μία για την ενεργό ισχύ $P$ και μία για την άεργο ισχύ $Q$, οι οποίες προστέθηκαν γύρω από το σημείο λειτουργίας κάθε φορτίου. Το πλάτος των διακυμάνσεων ορίστηκε ίσο με $0.1\%$ της αντίστοιχης βασικής τιμής του φορτίου, ενώ ο λευκός θόρυβος φιλτραρίστηκε με συχνότητα αποκοπής $5$~Hz. Η διάρκεια της προσομοίωσης ορίστηκε στα $600$~s, με χρονικό βήμα ολοκλήρωσης $0.01$~s. Από την προσομοίωση εξήχθησαν οι μεταβλητές τάσης ($V$) και ρεύματος ($I$) για καθεμία από τις γεννήτριες του συστήματος.

\section{Προεπεξεργασία και N4SID}
\label{amb:preprocessing}

 Στόχος της διαδικασίας ταυτοποίησης είναι η εκτίμηση των ηλεκτρομηχανικών ρυθμών του συστήματος από τις μετρήσεις ambient αποκρίσεων. Για τον σκοπό αυτό εφαρμόζεται ο αλγόριθμος ταυτοποίησης χώρου καταστάσεων \textit{N4SID}~\cite{n4sid1994}, ο οποίος εκτιμά ένα διακριτό μοντέλο χώρου καταστάσεων και τις αντίστοιχες συχνότητες και αποσβέσεις των ρυθμών.

  Πριν από την εφαρμογή του αλγορίθμου, τα πρωτογενή δεδομένα υποβλήθηκαν σε \textit{detrending}, υποδειγματοληψία στα $5$~Hz και βαθυπερατό φιλτράρισμα με συχνότητα αποκοπής $2$~Hz. Ο \textit{N4SID} εφαρμόζεται επαναληπτικά σε δύο ομάδες τάξεων μοντέλου:
  \[
  n \in \{2,4,6,\dots,30\},
  \qquad
  n \in \{10,15,20,\dots,45\}.
  \]
  Μετά την ταυτοποίηση διατηρούνται μόνο οι εκτιμήσεις με συχνότητα στο διάστημα $[0.1,\,2.0]$~Hz και αρνητικό πραγματικό μέρος της ιδιοτιμής, δηλαδή ευσταθείς εκτιμήσεις, σύμφωνα με το εύρος των ηλεκτρομηχανικών ρυθμών του συστήματος~\cite{sarkar2024drga}.

\section{Ιδιοτιμές Αναφοράς από το PowerFactory}

Εκτός από τις χρονοσειρές τάσης και ρεύματος, από το \textit{PowerFactory} εξάγονται οι ιδιοτιμές αναφοράς για το ίδιο \textit{ambient} σενάριο. Η διαδικασία περιλαμβάνει \textit{modal analysis}, από το οποίο διατηρούνται οι ιδιοτιμές που σχετίζονται με ηλεκτρομηχανικές ταλαντώσεις. Ο διαχωρισμός βασίζεται στη συμμετοχή καταστάσεων των σύγχρονων γεννητριών, με έμφαση στη γωνία δρομέα ($\delta$) και στην ταχύτητα περιστροφής ($\omega$).

Για κάθε ρυθμό καταγράφονται η συχνότητα, η απόσβεση, οι γεννήτριες που συμμετέχουν και οι αντίστοιχες περιοχές ελέγχου. Διατηρούνται μόνο οι τρόποι με αρνητικό πραγματικό μέρος και θετικό φανταστικό μέρος. Το τελικό σύνολο αποτελεί την αναφορά για την αξιολόγηση των εκτιμήσεων N4SID και των αποτελεσμάτων της συσταδοποίησης.

\begin{table}[htbp]
    \centering
    \caption{Ηλεκτρομηχανικοί τρόποι αναφοράς που διατηρήθηκαν από το \textit{PowerFactory}.}
    \label{tab:amb:pf_reference_modes}
    \scriptsize
    \begin{tabular}{lccccc}
        \toprule
        \textbf{Mode} & \textbf{$f$ [Hz]} & \textbf{$\sigma$} & \textbf{Generators} & \textbf{Areas} & \textbf{Type} \\
        \midrule
        Mode 1 & 1.110 & -0.193 & 10, 8 & 1 & Intra-area \\
        Mode 2 & 0.611 & -0.385 & 1, 5, 6, 9 & 1--3 & Inter-area \\
        Mode 3 & 1.049 & -0.399 & 1--7, 10 & 1--3 & Inter-area \\
        Mode 4 & 0.968 & -0.464 & 5, 9 & 1, 3 & Inter-area \\
        Mode 5 & 1.196 & -0.465 & 2, 3 & 2 & Intra-area \\
        Mode 6 & 1.187 & -0.586 & 4--7 & 3 & Intra-area \\
        Mode 7 & 1.423 & -0.696 & 6, 7 & 3 & Intra-area \\
        Mode 8 & 1.410 & -0.710 & 8, 10 & 1 & Intra-area \\
        Mode 9 & 1.451 & -0.817 & 4, 5 & 3 & Intra-area \\
        \bottomrule
    \end{tabular}
\end{table}
 
\section{Μέθοδοι Συσταδοποίησης και Υπερπαράμετροι}

Το screened σύνολο εκτιμήσεων ομαδοποιείται ανά περιοχή ελέγχου με τις μεθόδους \textit{k-Means}~\cite{macqueen1967some}, \textit{k-Medoids}~\cite{kaufman2009finding}, \textit{DBSCAN}~\cite{dbscan1996}, \textit{OPTICS}~\cite{optics1999}, \textit{HDBSCAN}~\cite{hdbscan2015}, \textit{Gaussian Mixture Models} (GMM)~\cite{gmm2002} και Agglomerative Hierarchical Clustering (AHC)~\cite{agglomerative1963}. Οι παράμετροι κάθε μεθόδου και η πολιτική επιλογής τους παρουσιάζονται σε επόμενη ενότητα μαζί με τα αποτελέσματα, ώστε η σύγκριση να βασίζεται σε κοινό input dataset και σαφώς ορισμένες ρυθμίσεις.

\subsection{Διαδικασία Ρύθμισης Υπερπαραμέτρων}
\label{amb:clustering-tuning}
Η ρύθμιση των υπερπαραμέτρων των αλγορίθμων συσταδοποίησης πραγματοποιήθηκε ανεξάρτητα για κάθε συνδυασμό order sweep και περιοχής ελέγχου,
χρησιμοποιώντας το αντίστοιχο screened σύνολο εκτιμήσεων που προέκυψε απο τον N4SID. Για τις μεθόδους που απαιτούν εκ των προτέρων
καθορισμό αριθμού συστάδων ((\textit{k-Means}, \textit{k-Medoids} και \textit{Gaussian Mixture Models}) ή κατωφλίου γειτνίασης (\textit{DBSCAN}, \textit{HDBSCAN} και Agglomerative Hierarchical Clustering), εξετάστηκε ένα προκαθορισμένο σύνολο υποψήφιων τιμών. Η τελική επιλογή των υπερπαραμέτρων έγινε με βάση την μεγιστοποίηση του δείκτη \textit{silhouette score} για όλους τους αλγορίθμους συσταδοποίησης, με μοναδική εξαίρεση τον GMM για τον οποίο επιλέχθηκε η παραμετροποίηση που ελαχιστοποιεί το κριτήριο BIC.[Να γίνει ναφορά στο θεωρητικό υπόβαθρο που θα εξηγει τους 2 δείκτες αξιολόγησης]

Η συσταδοποίηση εφαρμόστηκε στις συντεταγμένες των πόλων
\[
\mathbf{X}=[\sigma,\omega],
\qquad
\omega=2\pi f,
\]
όπου η απόσβεση $\sigma$ και η γωνιακή συχνότητα $\omega$ εκφράζονται
σε rad/s. Οι ίδιες συντεταγμένες χρησιμοποιήθηκαν για την εφαρμογή των
αλγορίθμων, τον υπολογισμό των αποστάσεων και την αξιολόγηση των
υποψήφιων υπερπαραμέτρων. Τα modal maps παρουσιάζονται στο ισοδύναμο επίπεδο
$(\sigma,f)$, ώστε η συχνότητα να εκφράζεται σε Hz για λόγους αναγνωσιμότητας.

Για τις μεθόδους πυκνότητας (\textit{DBSCAN}, \textit{OPTICS} και
\textit{HDBSCAN}), οι παράμετροι $\epsilon$ και $N_{\mathrm{pts}}$
υπολογίστηκαν σύμφωνα με τη διαδικασία της
\cite{sfetkos2024measurement}:
\[
\epsilon =
\frac{p_e}{2}
\sqrt{(\sigma_{\max}-\sigma_{\min})^2+
    (\omega_{\max}-\omega_{\min})^2},
\]
\[
N_{\mathrm{pts}}=
\max\left(2,\left\lceil p_m N_s|n|\right\rceil\right),
\]
όπου $N_s$ είναι το πλήθος των συνδυασμών γεννήτριας$-$σήματος στην
εξεταζόμενη περιοχή ελέγχου και $|n|$ το πλήθος των τάξεων μοντέλου του
αντίστοιχου sweep. Εξετάστηκαν οι τιμές
$p_e \in \{0.01,0.015,\dots,0.15\}$ και
$p_m \in \{0.10,0.15,\dots,0.40\}$.

Για τον DBSCAN εξετάστηκαν όλοι οι συνδυασμοί $\epsilon$ και
$N_{\mathrm{pts}}$. Στον OPTICS το $N_{\mathrm{pts}}$ χρησιμοποιήθηκε
ως \texttt{min\_samples}, ενώ η παράμετρος \texttt{xi}
μεταβλήθηκε στο
διάστημα $\texttt{xi} \in \{0.02,0.04,\dots,0.40\}$. Στον HDBSCAN εξετάστηκαν οι τιμές
$N_{\mathrm{pts}}$ για την \texttt{min\_cluster\_size}, η τιμή
$\epsilon$ για την \texttt{cluster\_selection\_epsilon} και οι πολιτικές
\texttt{eom} και \texttt{leaf}. Στην Agglomerative Hierarchical
Clustering εξετάστηκαν η απόσταση κατωφλίου $\epsilon$ και οι μέθοδοι
σύνδεσης \textit{average} και \textit{complete}.

Για τις μεθόδους \textit{k-Means} και \textit{k-Medoids}, ο αριθμός
συστάδων μεταβλήθηκε στο διάστημα
\[
K=2,\dots,\min(10,N-1),
\]
όπου $N$ είναι το πλήθος των διαθέσιμων screened εκτιμήσεων. Για τα
Gaussian Mixture Models εξετάστηκαν
\[
K=1,\dots,\min(10,N),
\]
με πλήρη πίνακα συνδιακύμανσης, \texttt{reg\_covar}$=10^{-4}$,
\texttt{n\_init}$=10$ και αρχικοποίηση \textit{k-means++}.

Μία υποψήφια ρύθμιση θεωρήθηκε έγκυρη για την αξιολόγηση με silhouette
μόνο όταν παρήγαγε τουλάχιστον δύο συστάδες μη θορύβου. Δεν εφαρμόστηκε
πρόσθετο κριτήριο ελάχιστης κάλυψης των εκτιμήσεων.




\section{Μετρικές Αξιολόγησης}

Για την αξιολόγηση των αποτελεσμάτων, κάθε εκτιμώμενος ρυθμός αντιστοιχίζεται αρχικά στον πλησιέστερο σχετικό ηλεκτρομηχανικό ρυθμό αναφοράς του \textit{PowerFactory}. Η αντιστοίχιση αυτή χρησιμοποιείται ως βάση σύγκρισης με τις συστάδες που προκύπτουν από τις μεθόδους συσταδοποίησης.

Οι μετρικές \textit{V-measure} και ARI αξιολογούν τον βαθμό συμφωνίας μεταξύ
των συστάδων που προκύπτουν από κάθε μέθοδο και της παραπάνω αντιστοίχισης
στους ρυθμούς αναφοράς. Επομένως, εκφράζουν κατά πόσο η μέθοδος διαχωρίζει
τις εκτιμήσεις που αντιστοιχούν σε διαφορετικούς ηλεκτρομηχανικούς ρυθμούς.

Συμπληρωματικά, ο MAD υπολογίζεται για κάθε μέθοδο και για κάθε ρυθμό
αναφοράς. Εκφράζει τη διάμεση απόκλιση των εκτιμήσεων που αποδίδονται στον
συγκεκριμένο ρυθμό από την αντίστοιχη ιδιοτιμή αναφοράς, ως προς το
πραγματικό και το φανταστικό μέρος της ιδιοτιμής, δηλαδή στο επίπεδο
$(\sigma,\omega)$. Μικρότερες τιμές MAD υποδηλώνουν ακριβέστερη εκτίμηση του
αντίστοιχου ρυθμού.
